% figures/w05-two-integrators.svg — 같은 조준식, 두 가지 적분기.
%
%   bash _tools/tikz2svg.sh figures/src/w05-two-integrators.tex figures/w05-two-integrators.svg
%
% 사용자 지시 (2026-10-13): 2025년 강의자료 "2. USV 제어기 설계" 슬라이드 66-69 의
% 비교를 이 강의에 맞춰 다시 그린다. 원본은 Fossen 의 TTK 4190 강의노트 그림이므로
% **옮겨 붙이지 않고 같은 구조를 직접 그린다** (CLAUDE.md §2).
%
% 이 그림이 말하는 것 하나:
%   두 법칙은 **조준식이 같다**. 다른 것은 적분 상태를 갱신하는 한 줄뿐이고,
%   그 한 줄의 차이는 아래 그래프에 전부 들어 있다 — 비선형 쪽은 y_e = Delta 에서
%   1/2 로 머리를 치고 내려오고, 보통 적분은 경로에서 멀수록 끝없이 커진다.
%
% 배치 (mm): 공유 상자 (95,96) / 왼쪽 열 x=48, 오른쪽 열 x=146 / 그래프 (18,-34)
% 겹침을 피하려고 열 간격을 98 mm 로 벌렸다. 좁히면 수식 상자가 서로 물린다.
\documentclass[border=6pt]{standalone}
% gnc-style.tex — 이 볼트의 TikZ 그림이 공유하는 스타일.
%
% 저널 그림 규약을 스타일 하나로 굳혀 둔다. 그림마다 선 굵기나 화살촉을
% 다시 정하지 않는다 — 그것이 그림이 제각각으로 보이는 원인이다.
%
% 쓰는 법:  % gnc-style.tex — 이 볼트의 TikZ 그림이 공유하는 스타일.
%
% 저널 그림 규약을 스타일 하나로 굳혀 둔다. 그림마다 선 굵기나 화살촉을
% 다시 정하지 않는다 — 그것이 그림이 제각각으로 보이는 원인이다.
%
% 쓰는 법:  % gnc-style.tex — 이 볼트의 TikZ 그림이 공유하는 스타일.
%
% 저널 그림 규약을 스타일 하나로 굳혀 둔다. 그림마다 선 굵기나 화살촉을
% 다시 정하지 않는다 — 그것이 그림이 제각각으로 보이는 원인이다.
%
% 쓰는 법:  \input{gnc-style}  (figures/src/ 안에서)

\usepackage{tikz}
\usepackage{amsmath}
\usepackage{xcolor}
\usetikzlibrary{arrows.meta, positioning, calc, fit, backgrounds, decorations.pathmorphing}

% ---- 색 --------------------------------------------------------------------
% 색은 강조 하나에만 쓴다. 나머지는 검정.
\definecolor{ink}{HTML}{1A1A1A}
\definecolor{accent}{HTML}{7C3AED}
\definecolor{warn}{HTML}{B91C1C}
\definecolor{soft}{HTML}{666666}

% 기하 그림(좌표계·유도)에서만 쓰는 넷. 블록선도는 위의 흑백 + accent 로 충분하다.
\definecolor{pathc}{HTML}{0D9488}   % 계획 경로 · 항적
\definecolor{errc}{HTML}{C2410C}    % 오차 (cross-track 등)
\definecolor{alng}{HTML}{B45309}    % along-track
\definecolor{flow}{HTML}{0369A1}    % 조류 · 유속

% ---- 선과 화살촉 -----------------------------------------------------------
\tikzset{
  % 신호선. 화살촉은 작고 뾰족한 것 하나뿐이다.
  sig/.style   = {ink, line width=0.5pt, -{Stealth[length=4pt, width=3pt]}},
  wire/.style  = {ink, line width=0.5pt},
  % 강조 경로 — 파선으로 구분한다. 흑백 인쇄에서도 살아야 하기 때문이다.
  asig/.style  = {accent, line width=0.5pt, densely dashed,
                  -{Stealth[length=4pt, width=3pt]}},
  awire/.style = {accent, line width=0.5pt, densely dashed},
  % 블록. 흰 바탕, 직각, 검은 테두리.
  blk/.style   = {draw=ink, line width=0.55pt, fill=white,
                  minimum width=17mm, minimum height=8mm, inner sep=2pt, font=\small},
  % 합산점
  sum/.style   = {draw=ink, line width=0.55pt, fill=white, circle,
                  inner sep=0pt, minimum size=4.6mm, font=\scriptsize},
  asum/.style  = {draw=accent, line width=0.55pt, fill=white, circle,
                  inner sep=0pt, minimum size=4.6mm, font=\scriptsize},
  % 분기점
  dot/.style   = {ink, fill=ink, circle, inner sep=0pt, minimum size=1.5mm},
  % 신호 이름 — 선 위에 작게
  sname/.style = {font=\small, inner sep=1.5pt},
  % 그림 안의 짧은 설명. 그림 안의 글은 최소로 둔다
  note/.style  = {font=\scriptsize, soft, inner sep=1.5pt},
  anote/.style = {font=\scriptsize, accent, inner sep=1.5pt},
  % 축
  axis/.style  = {ink, line width=0.5pt},
  tick/.style  = {font=\scriptsize, soft},
}

% 캡션 한 줄 — 그림 아래 가는 선과 함께
\newcommand{\gnccaption}[2]{%
  \node[anchor=north west, text width=#1, align=left, font=\scriptsize, ink]
        at (current bounding box.south west) {\vspace{2pt}#2};}
  (figures/src/ 안에서)

\usepackage{tikz}
\usepackage{amsmath}
\usepackage{xcolor}
\usetikzlibrary{arrows.meta, positioning, calc, fit, backgrounds, decorations.pathmorphing}

% ---- 색 --------------------------------------------------------------------
% 색은 강조 하나에만 쓴다. 나머지는 검정.
\definecolor{ink}{HTML}{1A1A1A}
\definecolor{accent}{HTML}{7C3AED}
\definecolor{warn}{HTML}{B91C1C}
\definecolor{soft}{HTML}{666666}

% 기하 그림(좌표계·유도)에서만 쓰는 넷. 블록선도는 위의 흑백 + accent 로 충분하다.
\definecolor{pathc}{HTML}{0D9488}   % 계획 경로 · 항적
\definecolor{errc}{HTML}{C2410C}    % 오차 (cross-track 등)
\definecolor{alng}{HTML}{B45309}    % along-track
\definecolor{flow}{HTML}{0369A1}    % 조류 · 유속

% ---- 선과 화살촉 -----------------------------------------------------------
\tikzset{
  % 신호선. 화살촉은 작고 뾰족한 것 하나뿐이다.
  sig/.style   = {ink, line width=0.5pt, -{Stealth[length=4pt, width=3pt]}},
  wire/.style  = {ink, line width=0.5pt},
  % 강조 경로 — 파선으로 구분한다. 흑백 인쇄에서도 살아야 하기 때문이다.
  asig/.style  = {accent, line width=0.5pt, densely dashed,
                  -{Stealth[length=4pt, width=3pt]}},
  awire/.style = {accent, line width=0.5pt, densely dashed},
  % 블록. 흰 바탕, 직각, 검은 테두리.
  blk/.style   = {draw=ink, line width=0.55pt, fill=white,
                  minimum width=17mm, minimum height=8mm, inner sep=2pt, font=\small},
  % 합산점
  sum/.style   = {draw=ink, line width=0.55pt, fill=white, circle,
                  inner sep=0pt, minimum size=4.6mm, font=\scriptsize},
  asum/.style  = {draw=accent, line width=0.55pt, fill=white, circle,
                  inner sep=0pt, minimum size=4.6mm, font=\scriptsize},
  % 분기점
  dot/.style   = {ink, fill=ink, circle, inner sep=0pt, minimum size=1.5mm},
  % 신호 이름 — 선 위에 작게
  sname/.style = {font=\small, inner sep=1.5pt},
  % 그림 안의 짧은 설명. 그림 안의 글은 최소로 둔다
  note/.style  = {font=\scriptsize, soft, inner sep=1.5pt},
  anote/.style = {font=\scriptsize, accent, inner sep=1.5pt},
  % 축
  axis/.style  = {ink, line width=0.5pt},
  tick/.style  = {font=\scriptsize, soft},
}

% 캡션 한 줄 — 그림 아래 가는 선과 함께
\newcommand{\gnccaption}[2]{%
  \node[anchor=north west, text width=#1, align=left, font=\scriptsize, ink]
        at (current bounding box.south west) {\vspace{2pt}#2};}
  (figures/src/ 안에서)

\usepackage{tikz}
\usepackage{amsmath}
\usepackage{xcolor}
\usetikzlibrary{arrows.meta, positioning, calc, fit, backgrounds, decorations.pathmorphing}

% ---- 색 --------------------------------------------------------------------
% 색은 강조 하나에만 쓴다. 나머지는 검정.
\definecolor{ink}{HTML}{1A1A1A}
\definecolor{accent}{HTML}{7C3AED}
\definecolor{warn}{HTML}{B91C1C}
\definecolor{soft}{HTML}{666666}

% 기하 그림(좌표계·유도)에서만 쓰는 넷. 블록선도는 위의 흑백 + accent 로 충분하다.
\definecolor{pathc}{HTML}{0D9488}   % 계획 경로 · 항적
\definecolor{errc}{HTML}{C2410C}    % 오차 (cross-track 등)
\definecolor{alng}{HTML}{B45309}    % along-track
\definecolor{flow}{HTML}{0369A1}    % 조류 · 유속

% ---- 선과 화살촉 -----------------------------------------------------------
\tikzset{
  % 신호선. 화살촉은 작고 뾰족한 것 하나뿐이다.
  sig/.style   = {ink, line width=0.5pt, -{Stealth[length=4pt, width=3pt]}},
  wire/.style  = {ink, line width=0.5pt},
  % 강조 경로 — 파선으로 구분한다. 흑백 인쇄에서도 살아야 하기 때문이다.
  asig/.style  = {accent, line width=0.5pt, densely dashed,
                  -{Stealth[length=4pt, width=3pt]}},
  awire/.style = {accent, line width=0.5pt, densely dashed},
  % 블록. 흰 바탕, 직각, 검은 테두리.
  blk/.style   = {draw=ink, line width=0.55pt, fill=white,
                  minimum width=17mm, minimum height=8mm, inner sep=2pt, font=\small},
  % 합산점
  sum/.style   = {draw=ink, line width=0.55pt, fill=white, circle,
                  inner sep=0pt, minimum size=4.6mm, font=\scriptsize},
  asum/.style  = {draw=accent, line width=0.55pt, fill=white, circle,
                  inner sep=0pt, minimum size=4.6mm, font=\scriptsize},
  % 분기점
  dot/.style   = {ink, fill=ink, circle, inner sep=0pt, minimum size=1.5mm},
  % 신호 이름 — 선 위에 작게
  sname/.style = {font=\small, inner sep=1.5pt},
  % 그림 안의 짧은 설명. 그림 안의 글은 최소로 둔다
  note/.style  = {font=\scriptsize, soft, inner sep=1.5pt},
  anote/.style = {font=\scriptsize, accent, inner sep=1.5pt},
  % 축
  axis/.style  = {ink, line width=0.5pt},
  tick/.style  = {font=\scriptsize, soft},
}

% 캡션 한 줄 — 그림 아래 가는 선과 함께
\newcommand{\gnccaption}[2]{%
  \node[anchor=north west, text width=#1, align=left, font=\scriptsize, ink]
        at (current bounding box.south west) {\vspace{2pt}#2};}


\definecolor{knownc}{HTML}{1D4ED8}

\tikzset{
  shared/.style = {draw=accent, line width=1.1pt, fill=accent!5, rounded corners=2pt,
                   align=center, inner sep=6pt, font=\small},
  lawA/.style   = {draw=errc, line width=1.0pt, fill=errc!5, rounded corners=2pt,
                   align=center, inner sep=6pt, font=\small},
  lawB/.style   = {draw=knownc, line width=1.0pt, fill=knownc!5, rounded corners=2pt,
                   align=center, inner sep=6pt, font=\small},
  nte/.style    = {font=\scriptsize, soft},
  hdA/.style    = {font=\small\bfseries, errc},
  hdB/.style    = {font=\small\bfseries, knownc},
  capt/.style   = {font=\scriptsize\bfseries, warn},
  arr/.style    = {ink, line width=0.8pt, -{Stealth[length=4.4pt,width=3.0pt]}},
  axis2/.style  = {soft, line width=0.7pt, -{Stealth[length=4pt,width=3pt]}},
  curveA/.style = {errc, line width=1.4pt},
  curveB/.style = {knownc, line width=1.4pt},
  grid2/.style  = {soft!50, line width=0.3pt},
}

\begin{document}
\begin{tikzpicture}[x=1mm, y=1mm]

% =====================================================================
% 위 — 두 법칙이 **함께 쓰는** 조준식
% =====================================================================
\node[shared, text width=104mm] (SH) at (95,96)
  {\textbf{both laws aim the same way}\\[5pt]
   $\psi_d = \pi_p - \arctan\!\bigl(K_p\, y_e^{\,p} + K_i\, y_{\mathrm{int}}^{\,p}\bigr)$,
   \qquad $K_p = \dfrac{1}{\Delta}$, \quad $K_i = \kappa K_p$, \quad $\kappa > 0$};
\node[nte, anchor=north] at (95,82)
  {what differs is one line underneath: how the state $y_{\mathrm{int}}^{\,p}$ is updated};

\draw[arr] (48,78) -- (48,71);
\draw[arr] (146,78) -- (146,71);

% =====================================================================
% 왼쪽 — 비선형 적분 (Borhaug et al. 2008)
% =====================================================================
\node[hdA, anchor=north] at (48,71) {NONLINEAR ILOS \ (Børhaug et al.\ 2008)};
\node[lawA, text width=76mm, anchor=north] at (48,64)
  {$\dot y_{\mathrm{int}}^{\,p}
     = \dfrac{\Delta\, y_e^{\,p}}{\Delta^2 + \bigl(y_e^{\,p} + \kappa\, y_{\mathrm{int}}^{\,p}\bigr)^2}$};
\node[nte, anchor=north, text width=80mm, align=left] at (52,42)
  {The denominator grows with the error, so the state can never climb faster than
   $\tfrac12$ per second: \textbf{anti-windup written into the law} rather than bolted
   onto it. Børhaug et al.\ prove UGAS/ULES for it in cascade with a 3-DOF ship model.};

% =====================================================================
% 오른쪽 — 보통 적분 (교과서형)
% =====================================================================
\node[hdB, anchor=north] at (146,71) {CONVENTIONAL ILOS};
\node[lawB, text width=76mm, anchor=north] at (146,64)
  {$y_{\mathrm{int}}^{\,p} = \displaystyle\int_0^t y_e^{\,p}(\tau)\,\mathrm d\tau$
   \\[6pt] that is, \quad $\dot y_{\mathrm{int}}^{\,p} = y_e^{\,p}$};
\node[nte, anchor=north, text width=86mm, align=left] at (146,42)
  {The error enters as it stands, at the same rate twenty metres off the path as one
   metre off it. A long approach winds the state up, and the unwinding then carries
   the vessel through to the other side.};

% =====================================================================
% 아래 왼쪽 — 그 한 줄의 차이를 그래프 하나로 (y_int = 0 에서의 갱신 속도)
%   비선형:  D y / (D^2 + y^2)  ->  y = D 에서 최대 1/2
%   보통   :  y                  ->  끝없이 커진다
%   Delta = 5 m. 가로 0..20 m 를 60 mm, 세로 0..1 m/s 를 34 mm 로.
% =====================================================================
\begin{scope}[shift={(18,-34)}]
  \draw[grid2] (0,0) grid[xstep=15,ystep=8.5] (60,34);
  \draw[axis2] (0,0) -- (66,0);
  \draw[axis2] (0,0) -- (0,40);
  \node[nte, anchor=west] at (66.5,0) {$y_e^{\,p}$ \ [m]};
  %  축 이름표는 축 **왼쪽**에. 축 위에 두면 보통 적분의 직선과 겹친다
  \node[nte, anchor=south east, align=right] at (-0.5,41.5) {$\dot y_{\mathrm{int}}^{\,p}$ \ [m/s]};
  \foreach \v/\lx in {0/0, 5/15, 10/30, 15/45, 20/60}
    \node[nte, anchor=north] at (\lx,-1.2) {\v};
  \foreach \v/\ly in {0.5/17, 1.0/34}
    \node[nte, anchor=east] at (-1.2,\ly) {\v};

  %  보통 적분: 기울기 1. 화면 위로 바로 나간다
  \draw[curveB] plot[domain=0:1.18, samples=20] ({3*\x}, {34*\x});
  \draw[curveB, densely dashed] (3.54,40.1) -- (8,46);
  \node[knownc, font=\scriptsize, anchor=west] at (8.4,46)
       {$\dot y_{\mathrm{int}}^{\,p} = 20$ m/s at $y_e^{\,p} = 20$ m — forty times the cap};

  %  비선형: 5 y / (25 + y^2)
  \draw[curveA] plot[domain=0:20, samples=81]
        ({3*\x}, {34*(5*\x)/(25 + \x*\x)});
  \draw[errc, line width=0.5pt, densely dashed] (15,17) -- (0,17);
  \draw[errc, line width=0.5pt, densely dashed] (15,17) -- (15,0);
  \fill[errc] (15,17) circle (1.0);
  \node[errc, font=\scriptsize, anchor=west] at (17,20)
       {peak $\tfrac12$ at $y_e^{\,p} = \Delta$, then falling};

  \node[nte, anchor=north west, text width=66mm] at (0,-6)
       {The update rate of each law while the state is still small, at $\Delta = 5$ m.
        Both vanish only with $y_e^{\,p}$, which is what makes $y_e^{\,p} = 0$ the
        equilibrium in either case; only one of them is bounded on the way there.};
\end{scope}

% =====================================================================
% 아래 오른쪽 — 이 배에서 실제로 잰 것
% =====================================================================
\node[capt, anchor=north west, align=left, text width=76mm] at (108,4)
  {Measured on the Otter, Experiment 5-5c\\[4pt]
   \normalfont\scriptsize
   At the same $\kappa = 0.3$ the plain integral never settles: it circles the path,
   $11.5$ m of mean error over the last 100 s. Tuned down to $\kappa = 0.05$ it does
   follow, and still swings $13.5$ m on the way in, where the nonlinear law swings
   $3.9$ m and settles at $0.07$ m.\\[6pt]
   \normalfont\scriptsize
   The denominator is what lets $\kappa$ be chosen for the current instead of for the
   approach.};

\end{tikzpicture}
\end{document}
